% ============================================================================
% sec-intro.tex — Introduction, notation, and the shape of the results.
% Label prefix: int: (all labels owned by this file).
% ============================================================================
\section{Introduction, notation, and the shape of the results}\label{sec:intro}

Let $G$ be a finite connected simple graph on $n$ vertices. A \emph{path decomposition} of
$G$ is a partition of $E(G)$ into paths, and $\pn(G)$ denotes the minimum number of paths in
such a decomposition. In his 1966 Tihany lecture, published as \cite{Gallai68}, Gallai
conjectured that $\pn(G)\le\ceil{n/2}$ for every connected $G$; the conjecture was
communicated by Erd\H{o}s and is tracked as problem \#583 in \cite{EP583}. It remains open,
with $\floor{2n/3}$ the best general upper bound, due to Dean and Kouider \cite{DK00}.

The case analysis of the last six decades has been organized, to a remarkable degree, by the
structure of the subgraph $\EV(G)$ induced by the even-degree vertices, together with bounds
on the degrees of $G$ itself. Lov\'asz \cite{Lo68} showed that $E(G)$ decomposes into at most
$\floor{n/2}$ paths and cycles, and settled the case $|\EV(G)|\le1$; Pyber \cite{Pyb96}
proved that $\EV(G)$ a forest implies $\pn(G)\le\floor{n/2}$, with the general bound
$n/2+O(n^{3/4})$; Favaron and Kouider \cite{FK88} settled Eulerian graphs of maximum
degree~$4$; and Fan \cite{Fan05} proved the conjecture when every block of $\EV(G)$ is
triangle-free with maximum degree at most~$3$. Two strengthenings followed: Botler and
Sambinelli \cite{BS21} proved it for graphs whose $\EV(G)$ is a subgraph of a graph in their
class $\mathcal{G}$ (each block of $\EV(G)$ of maximum degree at most~$3$, at most one
triangular block per component; their main theorem is stated with Gallai's ceiling
$\ceil{n/2}$ rather than the floor of Theorem~\ref{int:thm-fan} below), and Fan, Hou and
Zhou \cite{FHZ23} proved it when $\EV(G)$ contains no pair of triangles and every block of
$\EV(G)$ has maximum degree at most~$3$. A parallel track bounds the degrees of $G$ itself:
$\Delta\le5$ \cite{BP19} and treewidth at most~$3$ \cite{BCLS20}. The class-based programme then accelerated: Zhang, Liu and Hong \cite{ZLH24}
proved the conjecture, with the sharper $\floor{n/2}$ bound, for all $3$-degenerate graphs
except $K_3$ and $K_5^-$; Lu and Yi \cite{LY24} did the same for treewidth at most~$4$ with
exceptions $K_3$, $K_5^-$, $K_5$; and Chen and Wu \cite{CW25} settled it for block graphs.
On the algorithmic side, exact path partitions on subcubic graphs are computable in
polynomial time \cite{MY26}. Two further general results bound the problem from above: for
triangle-free graphs, Chu, Fan and Zhou \cite{CFZ22} proved the global bound
$\pn(G)\le 3n/5$; for the \emph{odd semi-cliques} (Definition~\ref{int:def-osc}), the
question raised by Bonamy and Perrett \cite{BP19} whether $\pn(G)\le\floor{n/2}$ holds for
every connected graph that is not an odd semi-clique has its first substantial progress in
Chu, Fan and Zhou \cite{CFZ26}: if the E-subgraph ($\EV(G)$) is isomorphic to a complete
graph $K_m$ with $m\le15$, then $\pn(G)\le\floor{n/2}+1$. Finally, minimal path
decompositions of complete bipartite graphs were studied by Constantinou and Ellinas
\cite{CE18}, whose algorithms cover three of four parameter cases; we return to the open
case below.

\subsection{Blocks versus global maximum degree}\label{int:subsec-blockvsglobal}

The Fan-line conditions quoted above are \emph{block-level}: they constrain the blocks of
$\EV(G)$ one by one, not $\EV(G)$ as a whole. The distinction is not pedantic; the two
readings are not equivalent, in either direction. In $K_{1,5}$ with one pendant leaf attached
at each of its six vertices, the even-degree subgraph is $\EV(G)\cong K_{1,5}$, which has a
vertex of E-degree~$5$, while every block of $\EV(G)$ is a single edge (maximum degree~$1$):
a large \emph{global} degree can coexist with a Fan-line hypothesis satisfied block by block.
Conversely, a cut vertex of $\EV(G)$ may belong to arbitrarily many blocks each of maximum
degree at most~$3$, so the absence of high E-degree vertices is stronger than the block
condition. Accordingly, the correct reading of the twenty-year-old Fan-line hypothesis is:
\emph{every block of $\EV(G)$ is triangle-free and has maximum degree at most~$3$}
\cite{Fan05}, read through Theorem~2(c) of \cite{BS21}, with the forest case of
Theorem~2(b) of \cite{BS21} going back to Pyber \cite{Pyb96}. The theorems of
Section~\ref{subsec:astar} are new precisely at the block level: they admit a \emph{block} of
$\EV(G)$ of maximum degree at least~$4$.

\subsection{Contributions}\label{int:subsec-contrib}

The paper has four mathematical parts (Parts~I--III being constructive and Part~IV
conditional and machine-backed) and a verification apparatus. We list the contributions,
with pointers to the sections that carry them.

\begin{enumerate}[label=(\roman*),leftmargin=2em]
\item \textbf{The exact path number of $K_{4,n}$ for odd $n$, and a decision procedure for
the even columns of Case~4 of \cite{CE18}} (Section~\ref{subsec:k4odd}). Constantinou and
Ellinas \cite{CE18} partition the parameter pairs of $K_{n_1,n_2}$ into four cases and give
minimal-decomposition algorithms for three of them, leaving open Case~4 (odd $n_1$,
$1\le n_2\le n_1-1$) with an explicit request for its ``exact characteristics''. We close the
$n_2=4$ column completely:
\[
\pn(K_{4,n})=(n+1)/2 \qquad\text{for every odd } n\ge3,
\]
improving the best published bound for the family (the published bound $k+2$, where
$n=2k+1$, from \cite{Pyb96,Fan05,LY24}) by one. The value for odd $n\ge3$ is formalized in Lean as a
single theorem with axiom profile \texttt{[propext, Classical.choice, Quot.sound]} and no
\texttt{sorryAx}; numeric corollaries at $n=5,13,19$ follow from it without independent
gaps. For the remaining even columns we prove a general counting cap, a structure theorem
for optimal decompositions across the whole even range, and a \emph{ring-closure theorem}
(one explicit base decomposition of $K_{2k_0+1,2h}$ into $k_0+1$ paths extends, by an
explicit construction, to every $k\ge k_0+h+1$); the column $n_2=6$ closes completely, the
columns $n_2=8,10$ close except for five named base instances, and the conjecture that
$\pn(K_{n_1,n_2})=\ceil{n_1/2}$ throughout the even range (Conjecture~\ref{conj:con27A} in the
body) becomes a column-by-column decision procedure with $h+1$ base instances per column.

\item \textbf{Odd clique chains: refinements beyond block graphs, with the domain fixed}
(Section~\ref{subsec:chains}). Chen and Wu \cite{CW25} proved Gallai's conjecture for all block
graphs. For the subfamily of \emph{odd clique chains} --- block graphs whose block tree is a
path and whose blocks are $B_i\cong K_{2j_i+1}$ with \emph{all $j_i$ odd}, i.e.\ blocks of
order $\equiv 3 \pmod 4$, $K_3$ allowed --- we prove
\[
\pn(G)\le \floor{n/2}\ \ (b\ge2), \qquad \pn(G)\le \floor{n/2}-1\ \ (b\ge4),
\]
where $b$ is the number of blocks and $n=2\sum_i j_i+1$, along with the exact value
$\pn(G)=\ceil{n/2}$ for a single block. Two qualifications fix the domain precisely. First,
\emph{every} vertex of such a chain has even degree: a non-cut vertex of a block
$K_{2j_i+1}$ has degree $2j_i$, and a cut vertex shared by consecutive blocks has degree
$2j_i+2j_{i+1}$; hence $\EV(G)=G$. Second, the complementary case --- chains containing
blocks of order $1 \pmod 4$ (even $j_i$, e.g.\ $K_5$) --- is \emph{open}: no general
construction is known, and in-chain existence is machine-verified only for the block values
$K_5$ and $K_9$. For $b=3$ the $-1$ refinement is likewise open and reduces to a
distinctness question on the middle block. The local engine is: \emph{for every odd $j$ and
every Hamilton path $P$ of $K_{2j+1}$, the graph $K_{2j+1}-P$ decomposes into $j$ paths};
the base zigzag objects are classical \cite{LP92,Alspach08}, the statement for an
arbitrarily prescribed Hamilton path appears to be new.

\item \textbf{Fan-type theorems admitting a block of $\EV(G)$ of maximum degree $\ge4$}
(Section~\ref{subsec:astar}). The Fan line \cite{Fan05,BS21,FHZ23} has kept its requirement at
the block level for twenty years: every block of $\EV(G)$ has maximum degree at most~$3$.
We prove class theorems that admit a block of $\EV(G)$ of maximum degree at least~$4$:
Theorems~\ref{thm:claAstar} and~\ref{thm:lasep} (if $\EV(G)$ is triangle-free, high E-degrees are compensated by an
odd-degree condition, and an explicit alignment condition holds at one high vertex, then
$\pn(G)\le\floor{n/2}$; and the same bound with no alignment hypothesis
whenever distinct high vertices have no common EV-neighbour), and a two-high-vertex theorem
with a corner-case reduction. An honesty note travels with these results: our witness
family $W(m)$ ($m$ odd, $\ge5$) has $\EV(W(m))\cong K_{2,m}$, a single block of maximum
degree $m$, and has treewidth $5$, degeneracy $5$ and maximum degree at least $m+1$; it
therefore lies \emph{outside} the degeneracy and treewidth classes of \cite{ZLH24,LY24} and
outside every published class we are aware of, so our theorems and the class-based line
\cite{ZLH24,LY24,BCLS20,BP19,CW25} are incomparable rather than competitive. The general
statement is posed as Conjecture~\ref{conj:claAstar}.

\item \textbf{A conditional route: Gallai modulo an explicit premise ledger} (Part~IV,
Sections~\ref{sec:computation}--\ref{sec:open}). We develop an endpoint-budget framework: the
per-vertex \emph{endpoint budget} $\beta(G,x)$ (Definition~\ref{int:def-budget}) and the
\emph{budget conjecture} $K$ that $\beta(G,x)\le\ceil{n/2}$ always. On the proved side we
establish an absorption-cost viewpoint (Remark~\ref{rem:absorb}) unifying three obstruction
phenomena from the literature; that $K$ implies Gallai's conjecture
(Theorem~\ref{int:thm-Kgallai}, proved here), so $K$ is a genuine strengthening; and a
conditional chain (Theorem~\ref{int:thm-ledger}, stated here and carried at pointer level in
the ledger subsection) assembling
$\text{Gallai}\Leftarrow K\Leftarrow\text{degree-two layer of }K\Leftarrow
\text{SC-PH}\Leftarrow\text{a placement conjecture}\Leftarrow\text{(E)+(R)}$, with the
bridge from the valuation line (the criterion of Section~\ref{subsubsec:criteria}) carried as a
parallel, redundant-for-implication route. The chapter's central deliverable is an \emph{explicit
hypothesis ledger}: an eight-item premise set, restated descriptively in
Section~\ref{int:subsec-conditional} below. The premise set mixes three kinds of conjuncts:
external steps flagged as unverified against the journal record (the all-odd support of the
Lov\'asz decomposition theorem as consumed by our proof of
Theorem~\ref{int:thm-Kgallai}), open conjectures internal to the paper (BP-odd, the
odd-order density conjecture, SC-PH), and paper-open items with machine backing. No premise
is claimed to be proved.

\item \textbf{Verification methodology} (Sections~\ref{sec:computation} and the formalization
section, Appendices). All exact values and all decompositions cited in this paper were produced and
re-checked by an independent machine pipeline (a bitmask dynamic program plus a
decomposition verifier asserting that the parts are simple paths and exactly partition the
edge set). The large-scale machine accounts of Part~IV are quoted at their measured
strength, under two standing conventions. First, the \emph{one-way superset convention}:
every machine check of the covering identity is performed in a reading that omits one
definitional branch (the \texttt{v7} branch, whose machine face is frozen because the
legacy candidate-generation source files are lost), giving a conservative superset
$\mathrm{fam}_{\mathrm{noV7}}\subseteq\mathrm{Im}\cap L$; zero counterexamples on the
superset face strictly imply zero counterexamples on the true face, and not conversely (the
residual \texttt{v7}-only mass is nonzero and precisely measured on the tested domains).
Second, the \emph{accounting layers are kept separate}: the nine-domain key tables
(comprising $2{,}412{,}436$ keys and $98{,}602{,}396$ key-weighted positions over
$N\in\{6,\dots,14\}$) are counted key-weighted, each key once, while the per-instance
accounts at $N=8$ and $N=9$ are instance-weighted; the two accountings are never mixed, and
the $N=14$ domain is an affine-normalized domain whose equivalence to exhaustive enumeration
is unproven. Lean certificates carry their trust boundaries verbatim: theorems proved from
first principles are reported with their axiom profiles, and finite certificates checked by
\texttt{native\_decide} carry the kernel-reflection axiom \texttt{Lean.ofReduceBool}
explicitly.
\end{enumerate}

\subsection{Claim-strength declaration}\label{int:subsec-claims}

Six no-claim constraints bind the whole account, including this introduction; we state them
verbatim. We do \emph{not} claim:
\begin{enumerate}[label=\textup{(\arabic*)},leftmargin=2em]
\item closure at $L\ge13$ (machine evidence at larger $N$ is extrapolation evidence, not
proof; only candidate-level status is claimed there);
\item general-$N$ paper closure of the covering identity (the general-$N$ existence half
stays open);
\item that the four-piece (E)+(R) landing is complete;
\item paper closure of the general-$L$ uniform-construction account (the
capacity-starved branch of the channel-or-menu dichotomy programme);
\item full closure of that dichotomy programme;
\item an unconditional theorem-form statement of the criterion of
Conjecture~\ref{conj:clacand}. The Lean theorem \texttt{thmA2663\_upgraded} is a
\emph{criterion-level bound}: under the explicit hypotheses that the target side has the
prescribed size and that $\mathrm{val}\le1$, it proves $g_{13}\le d_2+1$ with axiom profile
\texttt{[propext, Quot.sound]} and no \texttt{sorryAx}; the hypothesis $\mathrm{val}\le1$
itself is \emph{not} proved by the theorem, and a per-instance criterion verified over
finitely many $N$ is not a feasibility-existence proof for general $N$.
\end{enumerate}

\subsection{Notation and preliminaries}\label{int:subsec-notation}

All graphs are finite, simple and undirected. A \emph{path} has at least one edge and no
repeated vertex; a \emph{path decomposition} of $G$ is a set $D$ of paths whose edge sets
partition $E(G)$; $\pn(G)$ is the minimum size of a path decomposition. For a decomposition
$D$ and a vertex $x$ we write $D(x)$ for the number of paths of $D$ having $x$ as an
endpoint. A \emph{block} of $G$ is a maximal $2$-connected subgraph or a bridge (together
with its endpoints); every two blocks share at most one vertex, a \emph{cut vertex}. $G$ is
\emph{$k$-degenerate} if every subgraph has a vertex of degree at most $k$. The
\emph{even-degree subgraph} $\EV(G)$ is the subgraph induced by the vertices of even degree;
for $v\in V(G)$ we write $\dEV(v)$ for the degree of $v$ inside $\EV(G)$, the \emph{E-degree}
of $v$; a vertex with $\dEV(v)\ge4$ is called \emph{high}. Note that $\dEV(v)\ge4$ implies
in particular that $v$ has even degree in $G$.

We collect the three folklore lower bounds used throughout; proofs are included for
completeness.

\begin{lemma}[odd-degree lower bound; endpoint parity]\label{int:lem-LB1}
Let $D$ be a path decomposition of $G$. Then:
\begin{enumerate}[label=(\roman*),leftmargin=2em]
\item $2|D|\ge \mathrm{odd}(G)$, where $\mathrm{odd}(G)$ is the number of odd-degree
vertices; in particular $\pn(G)\ge \mathrm{odd}(G)/2$;
\item every odd-degree vertex is an endpoint of at least one path of $D$;
\item for every vertex $v$, $D(v)\equiv \deg_G(v)\pmod 2$.
\end{enumerate}
\end{lemma}

\begin{proof}
For a path $P$ and a vertex $v$, the degree of $v$ in $P$ is $0$, $1$ or $2$, and it is $1$
exactly when $v$ is an endpoint of $P$. Summing over $D$,
$\deg_G(v)=\sum_{P\in D}\deg_P(v)\equiv D(v)\pmod 2$, which gives (iii); (ii) is (iii) for
odd $\deg_G(v)$; and (i) follows by charging each odd-degree vertex to one path ending at
it, each path being charged at most twice.
\end{proof}

\begin{lemma}[edge-count lower bound]\label{int:lem-LB2}
If $G$ is connected with $n\ge2$ vertices and $m$ edges, then
$\pn(G)\ge \ceil{m/(n-1)}$.
\end{lemma}

\begin{proof}
Each path has at most $n-1$ edges, so $k$ paths cover at most $k(n-1)$ edges.
\end{proof}

\begin{lemma}[disconnected sum]\label{int:lem-sum}
If $G$ has components $G_1,\dots,G_c$ on $n_1,\dots,n_c$ vertices, then
$\sum_i \floor{n_i/2}\le \floor{n/2}$.
\end{lemma}

\begin{proof}
Immediate from $\sum_i \floor{n_i/2}\le \floor{\sum_i n_i /2}$ for integers.
\end{proof}

The three standard tools from the literature, stated in the form used later:

\begin{theorem}[Lov\'asz \cite{Lo68}]\label{int:thm-lovasz}
Every $n$-vertex graph (not necessarily connected) has a decomposition of its edge set into
at most $\floor{n/2}$ paths and cycles.
\end{theorem}

\begin{theorem}[Pyber \cite{Pyb96}]\label{int:thm-pyber}
If $\EV(G)$ is a forest, then $\pn(G)\le \floor{n/2}$.
\end{theorem}

\begin{theorem}[Fan \cite{Fan05}]\label{int:thm-fan}
If every block of $\EV(G)$ is triangle-free with maximum degree at most $3$, then
$\pn(G)\le \floor{n/2}$.
\end{theorem}

\begin{definition}[odd semi-clique; {\cite{BP19}}]\label{int:def-osc}
An \emph{odd semi-clique} is a graph obtained from the complete graph on $2k+1$ vertices by
deleting at most $k-1$ edges.
\end{definition}

\begin{observation}[density obstruction]\label{int:obs-density}
An odd semi-clique on $n=2k+1$ vertices has
$m\ge \binom{2k+1}{2}-(k-1)=2k^2+1$ edges, so by Lemma~\ref{int:lem-LB2}
$\pn\ge \ceil{(2k^2+1)/2k}=k+1=\ceil{n/2}$; in particular no odd semi-clique admits a
$\floor{n/2}$-path decomposition.
\end{observation}

Finally we record the classical Walecki partition of the complete graph of odd order, in
the form used in Section~\ref{subsec:chains}; the base objects are classical
\cite{LP92,Alspach08}. Fix $j\ge 1$ and let $V=\{\infty\}\cup \mathbb{Z}_{2j}$. For
$k\in\mathbb{Z}_{2j}$ let $C_k$ be the cycle
\[
C_k \;=\; \infty,\, k,\, k+1,\, k-1,\, k+2,\, k-2,\;\dots,\, k+(j-1),\, k-(j-1),\, k+j,\, \infty
\]
(indices mod $2j$); equivalently the finite vertices of $C_k$ are $p_0=k$ and
$p_s=k+(-1)^{s+1}\ceil{s/2}$ for $s=1,\dots,2j-1$.

\begin{lemma}[Walecki partition; classical, cf.\ \cite{LP92,Alspach08}]\label{int:lem-walecki}
For $k=0,1,\dots,j-1$ the cycles $C_k$ are edge-disjoint and
$\bigsqcup_{k=0}^{j-1}E(C_k)=E(K_{2j+1})$.
\end{lemma}

\begin{proof}
Each $C_k$ is a cycle: the finite vertices $k, k\pm1,\dots,k\pm(j-1),k+j$ are $2j$ distinct
elements covering $\mathbb{Z}_{2j}$. We show every edge lies in exactly one $C_k$, $k$ taken
mod $j$. Write $\delta(x,y)$ for the circular distance of $x,y\in\mathbb{Z}_{2j}$, i.e.\ the
minimum of $\pm(x-y)$ mod $2j$ in $\{0,\dots,j\}$.

Reading off consecutive differences in the sequence $p_s$ (as above,
$p_{s+1}-p_s=\pm(s+1)$), one checks that the finite edges of $C_k$ are exactly the pairs
$\{k+t,\,k-t\}$ for $t\in\mathbb{Z}_{2j}\setminus\{0,j\}$ and the pairs
$\{k+t,\,k-t+1\}$ for $t\in\mathbb{Z}_{2j}$, together with the two $\infty$-edges
$\{\infty,k\}$, $\{\infty,k+j\}$. (In each of the two forms an edge is produced twice, once
from each endpoint, so it is counted once; $t=0$ and $t=j$ would give loops in the first
form and are excluded.) Hence $C_k$ has, among finite edges, two edges of each circular
distance $1,\dots,j-1$ and one edge of distance $j$.

Now fix an edge $e=\{x,y\}$ of $K_{2j+1}$ and let $d=\delta(x,y)$.
\begin{itemize}[leftmargin=1.5em]
\item If $y=x+j$ (the diameter case, $d=j$): for $j$ odd, $e$ has the form
$\{k+t,k-t+1\}$, which forces $2k\equiv 2x+d-1\pmod{2j}$, a unique solution mod $j$ (the
right side is even); for $j$ even, $e$ has the form $\{k+t,k-t\}$ with $2t\equiv d=j$,
forcing $2k\equiv 2x+j\pmod{2j}$, again a unique solution mod $j$. So $e$ lies in exactly
one $C_k$.
\item If $d<j$ is even: $e=\{x,x+d\}=\{k+t,k-t\}$ forces $2k\equiv 2x+d\pmod{2j}$, unique
mod $j$.
\item If $d<j$ is odd: $e=\{x,x+d\}=\{k+t,k-t+1\}$ forces $2k\equiv 2x+d-1\pmod{2j}$,
unique mod $j$.
\item If $e=\{\infty,x\}$: $e\in C_k$ iff $x\in\{k,k+j\}$, i.e.\ $k\equiv x$ or $x+j$
(mod $2j$), a unique class mod $j$.
\end{itemize}
In each case $e$ lies in exactly one of $C_0,\dots,C_{j-1}$, and each $C_k$ has $2j+1$
edges, so the $j(2j+1)=\binom{2j+1}{2}$ edges are partitioned.
\end{proof}

\subsection{Notation for the machine chapters}\label{int:subsec-machnotation}

The machine chapters (Section~\ref{sec:computation} and the formalization material of
Section~\ref{sec:open}) work with finite
combinatorial instances; we fix the notation once here, transcribed so that the chapter
statements are readable without further apparatus.

\begin{definition}[instances]\label{int:def-instance}
An \emph{instance} is $I=(N,P,T,d)$: $\mathbb{Z}_N:=\mathbb{Z}/N\mathbb{Z}$;
$P\subseteq\mathbb{Z}_N$ with $|P|=p$; $T=(t_1,t_2,t_3)$, three pairwise distinct elements
of $\mathbb{Z}_N$; $d=(d_1,d_2,d_3)$ with $d_i\ge1$; and the capacity surplus
$\beta:=p-d_1-d_2-d_3\ge1$ (capacity admission: $d_1\le p-[t_1\in P]$). The shifts are
$\delta_{12}:=t_1-t_2$, $\delta_{13}:=t_1-t_3$, $\delta_{23}:=t_2-t_3$ (mod $N$). The
instance is \emph{aligned} (\emph{so1}) when, with $g:=\gcd(\delta_{12},N)$ and
$L:=\mathrm{ord}(\delta_{12})$, one has $\delta_{13}\equiv k\,\delta_{12}\pmod N$ for the
canonical $k\in[2,L)$; equivalently, $t_1,t_2,t_3$ lie on one coset (``line'') of
$\langle\delta_{12}\rangle$.
\end{definition}

\begin{definition}[legal pairs and the value functional]\label{int:def-legalval}
For $C_1\subseteq P\setminus\{t_1\}$ put $B_{12}:=C_1-\delta_{12}$ and
$Q(C_1):=P\setminus(C_1\cup(B_{12}\cap P)\cup\{t_2\})$. A \emph{legal pair} is
$(C_1,C_2)$ with $C_1\subseteq P\setminus\{t_1\}$, $|C_1|=d_1$, $C_2\subseteq Q(C_1)$,
$|C_2|=d_2$. With $F_{13}:=(C_1-\delta_{13})\cap P$ and $F_{23}:=(C_2-\delta_{23})\cap P$,
the \emph{value} of a legal pair is the three-term count
\[
\mathrm{val}(C_1,C_2)
 :=\bigl|F_{13}\setminus(C_1\cup C_2)\bigr|
  +\bigl|F_{23}\setminus(C_1\cup C_2\cup F_{13})\bigr|
  +\bigl|\{t_3\}\cap P\setminus(C_1\cup C_2)\bigr| .
\]
The \emph{min-val} of the instance is the minimum of $\mathrm{val}$ over all legal pairs.
\end{definition}

The \emph{so1 (aligned) family} is the family of aligned instances; it is the family on
which the paper layer of the conditional chain is formulated, and the machine accounts of
Section~\ref{sec:computation} sweep exactly the statement ``every so1 instance with $\beta\ge2$
has $\min\mathrm{val}\le1$'' on the explicitly tested domains (the sweep domain $\beta\ge1$
contains the $\beta\ge2$ slice). The candidate form of that statement is registered as
Conjecture~\ref{conj:clacand} in Section~\ref{subsubsec:criteria}; its general-$N$ existence
half is Conjecture~\ref{int:conj-gapex}.

\subsection{Two reinsertion lemmas from the literature}\label{int:subsec-fanlemmas}

The Fan-type arguments of Section~\ref{subsec:astar} and the budget framework of
Section~\ref{int:subsec-conditional} consume two deletion--reinsertion lemmas, which we
state in the form used throughout; both follow the restatements in \cite{BS21}
(Lemmas~6 and~7 there), and we record our provenance boundaries with them.

\begin{lemma}[addible edge; {\cite{Fan05},\cite{BS21}}]\label{int:lem-fan6}
Let $uv\in E(G)$ and let $D'$ be a path decomposition of $G-uv$. If
$D'(v)>\ |\{w\in N_{G-uv}(u) : D'(w)=0\}|$, then $G$ has a path decomposition $D$ with
$|D|=|D'|$, $D(u)=D'(u)+1$, $D(v)=D'(v)-1$, and $D(w)=D'(w)$ for all other $w$. One says that
$uv$ is \emph{addible towards $u$}.
\end{lemma}

\begin{remark}[transcription, downgrade and machine evidence]\label{int:rem-fan6}
The proof of Lemma~\ref{int:lem-fan6} in \cite{Fan05} could not be consulted; the statement
follows the restatement in \cite{BS21} (Lemma~6) and has not been independently checked
against the original. A strengthening of the lemma that additionally requires the
underlying re-arrangement to preserve the endpoint profile of $D'$ pointwise is false in
general (a five-vertex counterexample is machine-certified); the lemma as stated imposes no
such requirement and has no known counterexample, and the Lean formalization used in
Section~\ref{sec:computation} carries the four-clause form obtained by deleting the final
pointwise clause $D(w)=D'(w)$ for the remaining vertices --- no application in this paper
consumes that clause. The machine evidence is direction-sensitive and has been
direction-corrected: the bidirectional rescan is complete at $n\le6$, all
$1{,}701{,}141{,}150$ ordered instances at $n=6$ (the fixed-orientation half exactly
reproducing the historical census of $850{,}570{,}575$), with zero violations
(Section~\ref{com:subsec-pipeline}).
\end{remark}

\begin{lemma}[star absorption; {\cite{Fan05},\cite{BS21}}]\label{int:lem-fan7}
Let $B=\{uw_1,\dots,uw_h\}$ be a set of $h$ edges incident to $u$, let $D'$ be a path
decomposition of $G-B$ in which every neighbour of $u$ \emph{in $G$} has $D'\ge 1$, and let
$uw$ be a specified edge of $B$. Then there is a set $A\subseteq B$ with $uw\in A$ and
$|A|\ge\ceil{h/2}$ such that the edges of $A$ can be added back one by one, each by
Lemma~\ref{int:lem-fan6} towards $u$, preserving $|D|$ and increasing $D(u)$ by one at each
step.
\end{lemma}

\begin{remark}\label{int:rem-fan7}
The hypothesis of Lemma~\ref{int:lem-fan7} is read over the neighbours of $u$ \emph{in
$G$}, so that it includes the star vertices $w_i$ themselves, which are not neighbours of
$u$ in $G-B$; this stronger reading is the one verified in every application below. As with
Lemma~\ref{int:lem-fan6}, the proof in \cite{Fan05} (Lemma~3.6 there) could not be
consulted, and the transcription from the restatement in \cite{BS21} (Lemma~7) has not been
independently checked against the original.
\end{remark}

\subsection{The absorption-cost accounting and the budget conjecture}\label{int:subsec-budget}

Several proofs in the literature --- Lov\'asz' deletion--reinsertion, Fan's star absorption,
and the analysis of minimal counterexamples --- fail at the same type of point: one needs a
decomposition of a reduced graph in which \emph{specified vertices hold sufficiently many
path endpoints}, and the count of available endpoints is not controlled by the usual
hypotheses. The following elementary accounting fact, quoted from the project's
verification work, quantifies the obstruction class once and is used only as a viewpoint in
this paper.

\begin{remark}[absorption cost]\label{rem:absorb}
Let $G$ be a graph, $D$ a path decomposition of $G$, $v$ a vertex not incident with any
edge of $G$ ($v$ may be a new vertex), and $N=\{w_1,\dots,w_d\}\subseteq V(G)$ a set of $d$
distinct vertices; let $G^+=G+\{vw: w\in N\}$. Build the bipartite graph $B(D)$ with left
side the $d$ edges $vw$ of the star and right side the paths of $D$, where $vw$ is adjacent
to $P$ iff $P$ has $w$ (or $v$) as an endpoint; let $\mu^*$ be a maximum matching of
$B(D)$. Within the operation class --- (i) attach a star edge $vw$ to the $w$-end of a path
of $D$ ending at $w$; (ii) pair two star edges $vw,vw'$ into the single path $w$--$v$--$w'$;
(iii) let a star edge be a single-edge path --- the minimum number of new paths required to
extend $D$ to a decomposition of $G^+$ is exactly $\ceil{(d-\mu^*)/2}$, and the resulting
$D(v)$ can be made at least $\mu^*$. The equality holds within the operation class and is
\emph{not} a lower bound for arbitrary decompositions of $G^+$. Three obstruction
phenomena in the literature are instances of ``the cost exceeds the budget'': the
Lov\'asz-style vertex reinsertion gap (the unproved strengthening of
Theorem~\ref{int:thm-lovasz} to path-only decompositions with endpoint control), the parity
wall for even-sized stars in the Fan engine (a star of even size leaves its centre even, so
no endpoint at the centre is forced), and the simplicity wall of Lemma~\ref{int:lem-fan6} (a
host endpoint may exist but be unusable because the path through it already contains the
reinsertion target).
\end{remark}

\subsection{The shape of the conditional route}\label{int:subsec-conditional}

The conditional chapters measure, at each vertex, the one quantity that the classical
deletion--reinsertion arguments leave uncontrolled: the number of path endpoints a
decomposition can be made to hold at a specified vertex.

\begin{definition}[endpoint budget]\label{int:def-budget}
For a connected graph $G$ and an even-degree non-isolated vertex $x$ of $G$, the
\emph{endpoint budget} $\beta(G,x)$ is the minimum cardinality of a path decomposition $D$
of $G$ in which $x$ is an endpoint of at least two paths, i.e.\
$\beta(G,x):=\min\{|D| : D \text{ a path decomposition of } G,\ D(x)\ge2\}$. The
\emph{budget conjecture} $K$ is the statement that $\beta(G,x)\le\ceil{n/2}$ for every
connected $G$ on $n$ vertices and every even-degree non-isolated vertex $x$.
\end{definition}

\begin{theorem}\label{int:thm-Kgallai}
Conjecture~$K$ implies Gallai's conjecture: every connected $n$-vertex graph $G$ has
$\pn(G)\le\ceil{n/2}$.
\end{theorem}

\begin{proof}
If $G$ has an even-degree non-isolated vertex $x$, then $K$ provides a decomposition with
at most $\ceil{n/2}$ paths and we are done. Otherwise $G$ is all-odd, so $n$ is even, and
Theorem~\ref{int:thm-lovasz} gives a decomposition into $p$ paths and $c$ cycles with
$p+c\le n/2$. In a mixed path-cycle decomposition the parity identity of
Lemma~\ref{int:lem-LB1}(iii) still holds (cycles contribute even degree at each vertex), so
every vertex is an endpoint of at least one path and $2p\ge n$; hence $p\ge n/2$, so
$p=n/2$, $c=0$, and the decomposition is pure, giving $\pn(G)\le n/2=\ceil{n/2}$.
\end{proof}

Thus $K$ is a strengthening of Gallai's conjecture, not a consequence of it. The remaining
links of the conditional chain --- from the degree-two layer of $K$ through the odd
semi-clique pair-hosting conjecture to the placement problems --- are assembled descriptively
below and carried at ledger level in Section~\ref{sta:subsec-limitations}; here we restate,
descriptively, the eight-item premise set behind the conditional-chain claim of
Section~\ref{int:subsec-contrib}; each link's verification level, premises and residual gaps
are registered at pointer level in the limitations subsection.

\begin{theorem}[conditional chain; descriptive registration]\label{int:thm-ledger}
There is a chain of implications, each available at paper level or registered at pointer
level with its verification status recorded, leading from the conjunction of the eight-item
premise set listed in this subsection (with the valuation criterion of
Conjecture~\ref{conj:clacand} entering, if at all, only through the bridge premises as a
parallel, redundant-for-implication route) to the conclusion
$\pn(G)\le\ceil{n/2}$ for every connected $G$ on $n$ vertices. The theorem is the
\emph{existence of the conditional route}, not the truth of any premise: every premise is an
open conjecture, an unverified external step, or a paper-open item with machine backing,
none is claimed proved, and no unconditional progress on Gallai's conjecture is claimed.
\end{theorem}

Each item below is an \emph{open} conjecture, an \emph{unverified external step}, or a
\emph{paper-open item with machine backing}; none is claimed proved.
\begin{itemize}[leftmargin=2em,itemsep=0.1em]
\item \emph{(external, flagged unverified against the journal record)} the all-odd
branch of the Lov\'asz decomposition theorem (Theorem~\ref{int:thm-lovasz}) as consumed by
the proof that the budget conjecture implies Gallai's conjecture: if $G$ is all-odd, the
theorem gives $p$ paths and $c$ cycles with $p+c\le n/2$, and the parity identity forces
$p=n/2$, $c=0$.
\item \emph{(external-origin, open)} BP-odd: every connected graph $H$ on an odd
number $2k+1$ of vertices with $|E(H)|\le 2k^2$ satisfies $\pn(H)\le k$ --- the odd-order,
tightness-special case of Question~1.1 of \cite{BP19}.
\item \emph{(internal, open)} the odd-order density conjecture at $k\ge4$: every
tight connected graph on $2k+1$ vertices ($\pn=\ceil{(2k+1)/2}$) is an odd semi-clique;
proved at $k\le3$ by exhaustive census, machine-certified at $k=4$ and supported by directed
search at $k=5$, open at $k\ge6$.
\item SC-PH: every odd semi-clique on $2k+1$ vertices hosts any two specified
non-isolated vertices on distinct paths of a $(k+1)$-path decomposition.
\item \emph{(open)} the surviving layers of the budget conjecture $K$: the
degree-two cut-vertex branch (computational level, awaiting theorization) and the
degree-$\ge4$ general-layer residual (a recorded dead zone of $42$ instances, discharged at
witness level only). \textbf{This premise is open as stated}: the recorded instances have
witness certificates, but no general paper proof exists, and the premise must be quoted as
an open hypothesis wherever the chain is used.
\item \emph{(paper-open, machine-backed)} general solvability of the placement
problem on the $(E)$-side, cyclic domain included: zero counterexamples at $k\le7$ over all
$188$ types (and machine accounts beyond), paper proof open.
\item \emph{(paper-open)} the $(R)$-side on the $3f\ge k$ domain: computational at
$k\le7$ (zero counterexamples), paper proof open.
\item \emph{(bridge, open)} the valuation-line bridge conditions: three registered
open conjectures of the difference-board analysis, quoted here descriptively and carried at
pointer level in the ledger.
\end{itemize}

We record explicitly what the ledger is and is not. The chain above is \emph{not} assembled
into a single master theorem of the form ``premises $\wedge$ criterion $\Rightarrow$ Gallai''
in this paper: the bridge premises (the last three items) are registered at pointer level
only, their full definitions residing in the project archive, so the paper certifies their
provenance rather than their re-adjudicability. What the ledger delivers is checkability:
any future resolution of any single premise propagates along the described chain through
paper-level implications whose verification status is recorded. No premise is claimed to be
within reach of the machine layer except as recorded, and the conditional route must not be
cited as progress towards Gallai's conjecture beyond its registered content.

\subsection{Reproducibility}\label{int:subsec-repro}

The machine pipeline behind every computational claim in this paper --- the exact-value
computations, the decomposition verifiers, the per-instance sweeps of
Section~\ref{sec:computation}, and the Lean development reported there and in
Section~\ref{sec:open} --- will be released together with the paper: the verification scripts, the certificate and
ledger files, the Lean sources with their axiom profiles, and the per-theorem verification
status table. The release makes every machine-checked number in this paper independently
re-runnable and every Lean theorem independently re-compilable.

% end sec-intro.tex
